% TOP_PROBLEM: {"id": 3196, "title": "Petersen coloring conjecture", "category_id": 3, "category": {"id": 3, "name": "graph_theory", "display_name": "Graph Theory", "description": "Problems involving graphs, networks, and their properties.", "slug": "graph-theory", "order_index": 3, "created_at": "2026-07-31T15:26:25.671Z"}, "status": "open", "problem_number": "OPG-143", "set_id": 12}
% FIRST_POSTED: 2026-10-01
% ENTRY_KIND: statement_only
% UnsolvedMath ID 3196; source catalogue code OPG-143
% !TeX program = lualatex
% Definition and sources only; this file is not an LLM solution attempt.
\documentclass[11pt,a4paper]{article}
\usepackage[margin=25mm]{geometry}
\usepackage{fontspec}
\setmainfont{lmroman10-regular.otf}[BoldFont=lmroman10-bold.otf,
 ItalicFont=lmroman10-italic.otf,BoldItalicFont=lmroman10-bolditalic.otf]
\usepackage{amsmath,amssymb,mathtools}
\usepackage{unicode-math}
\setmathfont{latinmodern-math.otf}
\AtBeginDocument{\let\setminus\smallsetminus}
\usepackage{xurl,hyperref}
\hypersetup{colorlinks=true,urlcolor=blue}
\setcounter{secnumdepth}{0}
\setlength{\parindent}{0pt}
\setlength{\parskip}{5pt}
\setlength{\emergencystretch}{3em}
\begin{document}
\section*{Petersen coloring conjecture}\label{3196}
{\small UnsolvedMath ID 3196 · OPG-143 · Graph Theory · OpenGarden}
\subsection{Definitions and mathematical statement}
Define the Petersen graph $P$ as the graph whose
vertices are the two-element subsets of
$\{1,2,3,4,5\}$, with two vertices adjacent
exactly when the subsets are disjoint. It has
ten vertices, each of degree three. For any
loopless graph $H$ and vertex $v$, let
$\delta_H(v)$ denote its set of incident edges.

Let $G$ be a finite loopless cubic multigraph
with no bridge. A Petersen colouring of $G$
is a function $f:E(G)\to E(P)$ satisfying
\[
 \forall v\in V(G)\quad\exists w\in V(P):
 \quad f\big|_{\delta_G(v)}:
 \delta_G(v)\longrightarrow\delta_P(w)
 \text{ is a bijection}.
\]
In particular, the three distinct edges at
each vertex of $G$ receive three distinct
Petersen-edge labels forming one complete
star of $P$. The same label may be used at
unrelated locations in $G$.

The conjecture asserts that every bridgeless
cubic $G$ admits such a map. The function
colours edges, not vertices; it is not required
to be an isomorphism, a covering map, or
surjective onto all fifteen edges of $P$.
The chosen vertex $w$ may depend on $v$.
The local star formulation makes the meaning
of the source's mutually adjacent edge triples
precise. Since the Petersen graph is triangle-free,
three distinct pairwise adjacent edges there
necessarily form a star. Every adjacent pair
of source edges receives distinct incident
labels by the displayed condition.
\subsection{Short English statement}
Can the edges of every bridgeless cubic graph be labelled by Petersen-graph edges so each vertex sees exactly the three edges at some Petersen vertex?
\subsection{Sources}
\begin{itemize}
\item[S1] ULAM.AI and the UnsolvedMath Contributors, supplied problems.json, record 3196 (OPG-143), OpenGarden collection. Catalogue identity and source transcription; CC BY 4.0.
\url{https://huggingface.co/datasets/ulamai/UnsolvedMath}.
\item[S2] Open Problem Garden, Petersen coloring conjecture. Original problem and discussion; the mathematical formulation below is expanded from the supplied source record.
\url{http://www.openproblemgarden.org/op/petersen_coloring_conjecture}.
\item[S3] F. Pirot, J.-S. Sereni and R. Škrekovski, Variations on the Petersen colouring conjecture, Electronic Journal of Combinatorics 27 (2020), P1.8; edge-map interpretation.
\url{https://arxiv.org/abs/1905.07913}.
\end{itemize}
\end{document}
